%% ============================================================
%%  REFERÊNCIA — pacote handout1colRL (coluna única)
%%  Copie este arquivo para começar um handout novo.
%%
%%  Compilar:  latexmk -xelatex referencia-1col.tex
%%             (ou xelatex duas vezes — o rodapé "n/total" precisa disso)
%%
%%  Troque a linha do \usepackage por
%%      \usepackage[notas]{handout1colRL}
%%  para abrir a margem lateral de anotações. Nada mais muda:
%%  \aomargem funciona nos dois modos.
%% ============================================================
\documentclass[a4paper, 11pt]{extarticle}

\usepackage[notas]{handout1colRL}
\usepackage{babel}
\babelprovide[import, main]{brazilian}

\disciplina{Aprendizagem por Reforço}
\rotulo{Folha de estudo}
\title{Aula N — Título da folha}
\subtitulo{Subtítulo curto \textbullet\ tópicos principais}
\author{Prof.\ Marcelo Keese Albertini \textbullet\ Universidade Federal de Uberlândia}
\rodape{Aprendizagem por Reforço \textbullet\ Aula N \textbullet\ UFU}

\begin{document}
\cabecalho

\secao{Caixas temáticas}

As mesmas oito caixas do deck Beamer, com as mesmas assinaturas — o
conteúdo pode ser copiado de um formato para o outro sem edição.
\aomargem{Notas laterais entram com \texttt{\textbackslash aomargem}.
Sem a opção \texttt{notas}, elas viram notas no corpo do texto.}

\begin{defbox}{Horizonte $H$}
  Número de passos de tempo em cada episódio.
\end{defbox}

\begin{defbox}{}
  Com título vazio, a caixa escreve apenas ``Definição''.
\end{defbox}

\begin{teobox}{Proposição — melhoria monotônica}
  $V^{\pol_{i+1}} \ge V^{\pol_i}$, com desigualdade estrita se $\pol_i$ é subótima.
\end{teobox}

\begin{algbox}{Iteração de Valor}
  \begin{itemize}
    \item Inicialize $V_0(s) = 0$ para todo $s$.
    \item Repita até convergir:
      $V_{k+1}(s) = \max_a\bigl[\rec(s,a) + \gamma\sum_{s'}\tran{s'}{s,a}V_k(s')\bigr]$.
  \end{itemize}
\end{algbox}

\begin{ideiabox}
  O valor de hoje se decompõe em: o que ganho agora mais o valor de onde posso parar.
\end{ideiabox}

\begin{alertabox}
  Busca exaustiva é inviável: há $\abs{\gA}^{\abs{\gS}}$ políticas determinísticas.
\end{alertabox}

\begin{exbox}{— Mars Rover}
  Calcule $V_{k+1}(s_6)$ com $\gamma = 0{,}9$.
\end{exbox}

\begin{respbox}
  Omita esta caixa na versão entregue aos alunos, se quiser que eles tentem antes.
\end{respbox}

\begin{quizbox}{}
  Um fator de desconto $\gamma$ grande dá mais peso a recompensas de curto prazo?
\end{quizbox}
\paradiscussao

\secao{Títulos e blocos}

\texttt{\textbackslash secao} desenha o filete laranja;
\texttt{\textbackslash subsecao} é a versão menor, sem filete. Para impedir
que um título fique órfão no pé da página, envolva título e conteúdo no
ambiente \texttt{juntos}.

\begin{juntos}
\subsecao{Bloco indivisível}
Este subtítulo e este parágrafo migram juntos para a página seguinte se não
couberem aqui.
\end{juntos}

\secao{Macros matemáticas}

As mesmas do deck: $\gS$, $\gA$, $\pol$, $\polstar$, $\Vpi$, $\Qpi$, $\Vstar$,
$\E$, $\Prob$, $\rec$, $\tran{s'}{s,a}$, $\abs{\gS}$, $\norminf{V_k - V_j}$,
$\defeq$, $\argmax_a Q(s,a)$.

Realce dentro de fórmula e anotação sob um termo:
\[
  \mdestaque{rlLaranja}{V = (I - \gamma P)^{-1}\rec}
  \qquad
  V(s) = \termo{\rec(s)}{recompensa imediata}
       + \termo{\gamma\textstyle\sum_{s'}\tran{s'}{s}V(s')}{futuro descontado}
\]

\secao{Diagramas}

\begin{center}
  \marsrover{4}
\end{center}
\notinha{\texttt{\textbackslash marsrover\{4\}} — destaca $s_4$.}

\begin{center}
  \marsrover*{}
\end{center}
\notinha{\texttt{\textbackslash marsrover*\{\}} — terminais em verde e rótulos
de recompensa.}

\begin{center}
  \scalebox{0.9}{\backupdiagrama}
\end{center}
\notinha{Para encolher um diagrama use \texttt{\textbackslash scalebox}, nunca
\texttt{scale=} — este último não reduz os nós.}

\secao{Tabelas}

Prefira \texttt{tabelaflex}, cujas colunas \texttt{Y} se ajustam sozinhas à
largura da página:

\begin{tabelaflex}{@{}p{0.26\linewidth}YYY@{}}
  \toprule
  & \colMC & \colTD & \colEC\\
  \midrule
  Precisa do modelo? & não & não & não (estima um)\\
  Faz bootstrapping? & não & sim & sim (via PD)\\
  Assume Markov?     & não & sim & sim\\
  \bottomrule
\end{tabelaflex}

\notinha{Existe também \texttt{tabelaRL}, com larguras fixas, para quando você
quiser controle exato. Ambas aceitam o corpo da fonte como argumento opcional:
\texttt{[\textbackslash scriptsize]}.}

\secao{Ajustes finos}

\begin{itemize}
  \item Margens: chame \texttt{\textbackslash geometry\{...\}} \emph{depois} de
        carregar o pacote.
  \item Cores pelo nome, com a mistura do xcolor (\texttt{rlAzul!15}):
        \texttt{rlAzul}, \texttt{rlAzulClaro}, \texttt{rlTeal},
        \texttt{rlLaranja}, \texttt{rlAmbar}, \texttt{rlVermelho},
        \texttt{rlVerde}, \texttt{rlCinza}, \texttt{rlFundo},
        \texttt{rlLinha}.
  \item Toda caixa aceita opções do tcolorbox no argumento opcional, como
        \texttt{[fontupper=\textbackslash small]}.
  \item O ambiente \texttt{colunas} existe como \emph{no-op}: um handout escrito
        para o \texttt{handoutRL} compila aqui em coluna única trocando só o
        nome do pacote.
\end{itemize}

\begin{teobox}[fontupper=\small]{Caixa com opção de tcolorbox}
  Este texto está menor porque a caixa recebeu \texttt{fontupper} no argumento
  opcional.
\end{teobox}

\end{document}
